\input{estilo}
% Cada chunk del archivo es el slice que lee un mapper: un archivo de n chunks,
% n mappers. La pizarra dibuja la flecha solo desde el primer chunk; acá sale
% una por chunk, y los puntos dicen que siguen hasta el n.
\begin{document}
\begin{figura}\begin{tikzpicture}
  \def\w{1.6}\def\h{0.72}
  \fill[relleno] (0,0) rectangle (\w,5*\h);
  \draw (0,0) rectangle (\w,5*\h);
  \foreach \k in {1,...,4} \draw (0,\k*\h) -- (\w,\k*\h);
  \foreach \k/\t [evaluate=\k as \y using (5-\k+0.5)*\h] in {1/1, 2/2, 3/3, 4/, 5/\var{n}} {
    \ifx\t\empty
      \node at (\w/2,\y) {$\vdots$}; \node at (\w+2.6,\y) {$\vdots$};
    \else
      \node[etiq] at (\w/2,\y) {chunk \t};
      \node[caja, minimum width=17mm, minimum height=5.5mm, font=\footnotesize] (m\k) at (\w+2.6,\y) {mapper \t};
      \draw[msg] (\w+0.1,\y) -- (m\k.west);
    \fi
  }
  \node[etiq, above=4pt] at (\w/2,5*\h) {archivo};
  \node[etiq, above=4pt] at (\w+1.3,5*\h) {slice = chunk};
\end{tikzpicture}\end{figura}
\end{document}
